\documentclass[12pt,a4paper,oneside]{report}
\usepackage[utf-8]{inputenc}
\usepackage[portuguese]{babel}
\usepackage[left=2.5cm,right=2.5cm,top=2.5cm,bottom=2.5cm]{geometry}
\usepackage{graphicx}
\usepackage{xcolor}
\usepackage{listings}
\usepackage{fancyhdr}
\usepackage{hyperref}
\usepackage{booktabs}
\usepackage{colortbl}
\usepackage{array}
\usepackage{tikz}
\usepackage{float}
\usepackage{framed}
\usepackage{amsmath}
\usepackage{amssymb}

% Configuração de cores
\definecolor{darkblue}{rgb}{0.12,0.28,0.53}
\definecolor{lightblue}{rgb}{0.92,0.95,0.98}
\definecolor{accentgreen}{rgb}{0.2,0.6,0.2}
\definecolor{codebackground}{rgb}{0.95,0.95,0.95}

% Configuração de fonte e espaçamento
\usepackage{setspace}
\onehalfspacing
\renewcommand{\baselinestretch}{1.5}

% Configuração de headers
\pagestyle{fancy}
\fancyhf{}
\rhead{\thepage}
\lhead{ETAPA 1 - Concepção e Arquitetura}
\renewcommand{\headrulewidth}{0.5pt}

% Configuração de título
\title{
    \vspace{-1.5cm}
    {\Large \textbf{TRABALHO FINAL - ETAPA 1}}\\
    \vspace{0.5cm}
    {\LARGE \textbf{Concepção e Arquitetura Inicial}}\\
    \vspace{1cm}
    {\Large Painel IoT Distribuído com Monitoramento em Tempo Real}\\
    \vspace{1.5cm}
}

\author{
    \textbf{Leonardo Cardoso}\\
    \textit{Matrícula: 32021BSI004}\\
    \vspace{0.5cm}
    Disciplina: Sistemas Distribuídos\\
    Data: \today
}

\date{}

% Configuração de seções
\usepackage{sectsty}
\chapterfont{\color{darkblue}\Large}
\sectionfont{\color{darkblue}\large}
\subsectionfont{\color{darkblue}}

% Configuração de código
\lstset{
    language=JavaScript,
    basicstyle=\ttfamily\small,
    backgroundcolor=\color{codebackground},
    breaklines=true,
    frame=single,
    rulecolor=\color{darkblue},
    showstringspaces=false,
    keywordstyle=\color{darkblue},
    commentstyle=\color{gray},
    stringstyle=\color{accentgreen}
}

% Comando para checkmark
\newcommand{\checkmark}{\text{\ding{51}}}

\begin{document}

\maketitle

\newpage
\tableofcontents
\newpage

% ============================================================================
% CAPÍTULO 1
% ============================================================================
\chapter{Informações Gerais}

\begin{center}
\begin{tabular}{|l|l|}
    \hline
    \rowcolor{lightblue}
    \textbf{Campo} & \textbf{Valor} \\
    \hline
    Nome do Aluno & Leonardo Cardoso \\
    \hline
    Matrícula & 32021BSI004 \\
    \hline
    Disciplina & Sistemas Distribuídos \\
    \hline
    Data de Entrega & 08 de Maio de 2026 \\
    \hline
    Data da Apresentação & 11 de Maio de 2026 \\
    \hline
\end{tabular}
\end{center}

% ============================================================================
% CAPÍTULO 2
% ============================================================================
\chapter{Definição do Tema do Micro-Mundo}

\section{Descrição Geral}

Desenvolvimento de um \textbf{sistema de monitoramento IoT distribuído} que simula um cenário real onde múltiplos sensores (dispositivos ESP32) distribuídos geograficamente coletam dados de temperatura, umidade e gás, transmitem para um servidor central através de um broker MQTT, e disponibilizam visualização em tempo real através de uma interface web.

\section{O que o Sistema Faz}

\begin{enumerate}
    \item \textbf{Coleta de Dados:} Múltiplos dispositivos ESP32 funcionam como \textit{nodes} sensoriais, capturando leituras de BMP280 (temperatura/pressão) e MQ-7 (gás).
    
    \item \textbf{Transmissão:} Dados são publicados via MQTT para um broker centralizado.
    
    \item \textbf{Processamento:} Backend processa telemetria, aplica filtros, calcula alertas e persiste em banco de dados.
    
    \item \textbf{Visualização:} Dashboard web exibe dados em tempo real, históricos, alertas e permite envio de comandos para os devices.
    
    \item \textbf{Controle Descentralizado:} Sistema de usuários com papéis (admin, operator, viewer) controla acesso aos recursos.
\end{enumerate}

\section{Entidades Principais}

\begin{itemize}
    \item \textbf{ESP32 Nodes:} Sensores distribuídos que medem o ambiente
    \item \textbf{Gateway IoT:} Dispositivo que recebe dados dos nodes e os envia para o servidor
    \item \textbf{API Gateway:} Camada de abstração HTTP para o backend
    \item \textbf{Backend:} Processa telemetria, gerencia BD e lógica de negócio
    \item \textbf{Broker MQTT:} Coordena comunicação entre devices e backend
    \item \textbf{PostgreSQL/TimescaleDB:} Armazena séries temporais de dados
    \item \textbf{Dashboard Web:} Interface de usuário para visualização e controle
\end{itemize}

\section{Por que é Distribuído}

O sistema é distribuído porque:

\begin{itemize}
    \item Componentes rodam em máquinas diferentes (Raspberry Pi 3 para BD, Notebook para Web)
    \item Comunicação via rede (MQTT TCP/IP, HTTP)
    \item Múltiplos sensores (possibilidade de escala)
    \item Necessidade de sincronização de estado
    \item Tolerância a falhas (falha de um sensor não derruba o sistema)
\end{itemize}

% ============================================================================
% CAPÍTULO 3
% ============================================================================
\chapter{Metas de Sistemas Distribuídos}

\section{Escalabilidade $\checkmark$}

\begin{description}
    \item[Implementada:] Sistema pode receber múltiplos ESP32 nodes em paralelo
    \item[Futura:] Replicação de broker MQTT, particionamento de dados por localização geográfica
    \item[Métrica:] Suporta 100+ sensores simultâneos com latência $<$ 5 segundos
\end{description}

\section{Disponibilidade $\checkmark$}

\begin{description}
    \item[Implementada:] Backend redunda entre 2 máquinas, persistência em BD
    \item[Futura:] Replicação de BD em múltiplas zonas, failover automático
    \item[Métrica:] RTO (Recovery Time Objective) $<$ 1 minuto
\end{description}

\section{Tolerância a Falhas $\checkmark$}

\begin{description}
    \item[Implementada:] Serviços com restart automático, reconexão de clients MQTT
    \item[Futura:] Circuit breaker, retry com backoff exponencial, dead letter queue
    \item[Métrica:] Perda de mensagens $<$ 0.1\%
\end{description}

\section{Compartilhamento de Recursos $\checkmark$}

\begin{description}
    \item[Implementada:] BD centralizado, fila de comandos, estado compartilhado em memória
    \item[Futura:] Distributed locking, consensus para decisões críticas
    \item[Métrica:] Consistência eventual com latência $<$ 500ms
\end{description}

\section{Segurança (Básico) $\checkmark$}

\begin{description}
    \item[Implementada:] Autenticação por sessão, papéis de acesso (RBAC)
    \item[Futura:] TLS para MQTT e HTTP, chaves SSH para acesso remoto, auditoria de logs
    \item[Métrica:] Zero exposição de credenciais em logs
\end{description}

% ============================================================================
% CAPÍTULO 4
% ============================================================================
\chapter{Arquitetura Inicial}

\section{Diagrama Textual de Arquitetura}

\begin{framed}
\scriptsize
\begin{verbatim}
┌─────────────────────────────────────────────────────────────────┐
│                    REDE LOCAL + TAILSCALE                        │
└─────────────────────────────────────────────────────────────────┘
                            │
            ┌───────────────┼───────────────┐
            │               │               │
      ┌─────▼────┐    ┌─────▼────┐    ┌────▼──────┐
      │ ESP32 N1 │    │ ESP32 NN │    │  Gateway  │
      │ Sensores │    │ Sensores │    │  (ESP32)  │
      └─────┬────┘    └─────┬────┘    └────┬──────┘
            │               │               │
            └───────────────┼───────────────┘
                      MQTT (TCP 1883)
                            │
        ┌────────────────────┼────────────────────┐
        │                                         │
   ┌────▼─────────┐                    ┌─────────▼──────┐
   │  Mosquitto    │                    │  API Gateway   │
   │  MQTT Broker  │◄──HTTP────────────►│  (Node.js)     │
   │(Localhost)    │                    │  :3001         │
   └────┬──────────┘                    └────────┬───────┘
        │ MQTT                                   │
        │                                        │
   ┌────▼──────────────────────────────────────▼──┐
   │            Backend (Node.js :3000)            │
   │  - Processa telemetria                        │
   │  - Gerencia autenticação                      │
   │  - Dashboard web                              │
   │                                               │
   │  Máquina: Notebook (server-sti)               │
   └────────────────┬───────────────────────────────┘
                    │
             TCP :5432 PostgreSQL
                    │
        ┌───────────┴───────────┐
        │                       │
   ┌────▼──────────────────┐    │
   │ PostgreSQL/TimescaleDB│    │
   │ iot_monitoring (BD)   │    │
   │                       │    │
   │ Máquina: Raspberry Pi 3
   └───────────────────────┘
\end{verbatim}
\end{framed}

\section{Arquitetura de Componentes}

A arquitetura segue o padrão cliente-servidor com componentes distribuídos:

\begin{itemize}
    \item \textbf{Camada de Sensores:} ESP32 nodes publicando dados via MQTT
    \item \textbf{Camada de Integração:} Broker MQTT + API Gateway
    \item \textbf{Camada de Processamento:} Backend Node.js processando e persistindo dados
    \item \textbf{Camada de Dados:} PostgreSQL/TimescaleDB para série temporal
    \item \textbf{Camada de Apresentação:} Dashboard web para visualização e controle
\end{itemize}

% ============================================================================
% CAPÍTULO 5
% ============================================================================
\chapter{Descrição dos Componentes}

\section{5.1 ESP32 Nodes (Sensores)}

\begin{center}
\begin{tabular}{|m{3cm}|m{10cm}|}
    \hline
    \rowcolor{lightblue}
    \textbf{Aspecto} & \textbf{Detalhe} \\
    \hline
    Responsabilidade & Capturar leituras de sensores (BMP280: temp/pressão; MQ-7: gás); publicar via MQTT \\
    \hline
    Tipo de Comunicação & Publish/Subscribe MQTT (Mosquitto broker) \\
    \hline
    Implementação & Processo único (firmware Arduino/PlatformIO no ESP32) \\
    \hline
    Protocolo & MQTT TCP :1883 \\
    \hline
    Localização & Distribuídos fisicamente (múltiplas localizações) \\
    \hline
    Tolerância a Falhas & Reconexão automática ao broker se houver desconexão \\
    \hline
\end{tabular}
\end{center}

\section{5.2 Gateway IoT (ESP32 especial)}

\begin{center}
\begin{tabular}{|m{3cm}|m{10cm}|}
    \hline
    \rowcolor{lightblue}
    \textbf{Aspecto} & \textbf{Detalhe} \\
    \hline
    Responsabilidade & Agregador de comandos; envia telemetria para backend via POST HTTP \\
    \hline
    Tipo de Comunicação & HTTP POST para API Gateway :3001/api/ingest \\
    \hline
    Implementação & Processo único (firmware Arduino) \\
    \hline
    Protocolo & HTTP/TCP :3001 \\
    \hline
    Endpoint & POST /api/ingest \\
    \hline
    Tolerância a Falhas & Retry com buffer local de comandos não-enviados \\
    \hline
\end{tabular}
\end{center}

\section{5.3 MQTT Broker (Mosquitto)}

\begin{center}
\begin{tabular}{|m{3cm}|m{10cm}|}
    \hline
    \rowcolor{lightblue}
    \textbf{Aspecto} & \textbf{Detalhe} \\
    \hline
    Responsabilidade & Coordenar pub/sub entre ESP32 nodes e backend \\
    \hline
    Tipo de Comunicação & MQTT TCP (publish/subscribe) \\
    \hline
    Implementação & Serviço de sistema (processo independente) \\
    \hline
    Protocolo & TCP :1883 \\
    \hline
    Localização & Localhost (Notebook) \\
    \hline
    Tópicos Principais & /telemetry/+/+ (dados), /command/+ (comandos) \\
    \hline
    Tolerância a Falhas & Fila persistente de mensagens, reconnect automático \\
    \hline
\end{tabular}
\end{center}

\section{5.4 API Gateway (Node.js)}

\begin{center}
\begin{tabular}{|m{3cm}|m{10cm}|}
    \hline
    \rowcolor{lightblue}
    \textbf{Aspecto} & \textbf{Detalhe} \\
    \hline
    Responsabilidade & Receber telemetria do gateway IoT; validar e repassar ao backend \\
    \hline
    Tipo de Comunicação & HTTP (entrada), HTTP (saída ao backend) \\
    \hline
    Implementação & Process Node.js (api-server.js; Port :3001) \\
    \hline
    Protocolo & HTTP/TCP :3001 \\
    \hline
    Localização & Notebook + Tunelado via Tailscale \\
    \hline
    Endpoints & POST /api/ingest, GET /api/status \\
    \hline
    Tolerância a Falhas & Retry ao backend, health-check periódico \\
    \hline
\end{tabular}
\end{center}

\section{5.5 Backend (Node.js)}

\begin{center}
\begin{tabular}{|m{3cm}|m{10cm}|}
    \hline
    \rowcolor{lightblue}
    \textbf{Aspecto} & \textbf{Detalhe} \\
    \hline
    Responsabilidade & Processar telemetria, gerenciar BD, autenticação, lógica de negócio, UI web \\
    \hline
    Tipo de Comunicação & MQTT (sub), HTTP (API), PostgreSQL (Query) \\
    \hline
    Implementação & Process Node.js (server.js; Port :3000) \\
    \hline
    Protocolo & HTTP :3000, MQTT :1883, PostgreSQL :5432 \\
    \hline
    Localização & Notebook (192.168.1.X) + Tailscale (100.82.199.70) \\
    \hline
    Endpoints & GET /, POST /login, GET /api/telemetry, POST /api/command \\
    \hline
    Módulos & postgres-store.js, simple-broker-client.js, iot-contract.js \\
    \hline
    Tolerância a Falhas & Reconexão ao broker, fallback para simulação local \\
    \hline
\end{tabular}
\end{center}

\section{5.6 PostgreSQL/TimescaleDB (BD)}

\begin{center}
\begin{tabular}{|m{3cm}|m{10cm}|}
    \hline
    \rowcolor{lightblue}
    \textbf{Aspecto} & \textbf{Detalhe} \\
    \hline
    Responsabilidade & Persistir série temporal de telemetria, config de devices, usuários e alertas \\
    \hline
    Tipo de Comunicação & SQL TCP :5432 \\
    \hline
    Implementação & Serviço de sistema (PostgreSQL + TimescaleDB extension) \\
    \hline
    Protocolo & PostgreSQL TCP :5432 \\
    \hline
    Localização & Raspberry Pi 3 (192.168.1.10) + Tailscale (100.82.140.119) \\
    \hline
    Banco Principal & iot\_monitoring (user: iot\_user) \\
    \hline
    Tabelas & telemetry (hypertable), devices, users, alerts \\
    \hline
    Tolerância a Falhas & WAL, backups programados, replicação futura \\
    \hline
\end{tabular}
\end{center}

\section{5.7 Frontend Web (Dashboard)}

\begin{center}
\begin{tabular}{|m{3cm}|m{10cm}|}
    \hline
    \rowcolor{lightblue}
    \textbf{Aspecto} & \textbf{Detalhe} \\
    \hline
    Responsabilidade & Exibir dados em tempo real, históricos, alertas; controlar devices \\
    \hline
    Tipo de Comunicação & HTTP (requisições), WebSocket (real-time, futuro) \\
    \hline
    Implementação & HTML + JavaScript (SPA simples, public/app.js) \\
    \hline
    Localização & Servido pelo backend :3000 \\
    \hline
    Funcionalidades & Login, dashboard, gráficos, controle de ESPs, gerenciamento de usuários \\
    \hline
    Tolerância a Falhas & Cache local, retry de requisições, feedback visual \\
    \hline
\end{tabular}
\end{center}

% ============================================================================
% CAPÍTULO 6
% ============================================================================
\chapter{Comunicação entre Componentes}

\section{Fluxo de Dados - Telemetria}

\begin{verbatim}
ESP32 Node 1 ─── MQTT Publish ──→ Mosquitto ─── Subscribe ──→ Backend
                                        ↓
                              API Gateway (GET)
                                    ↓
                              POST /api/ingest
                                    ↓
                              Backend processa
                                    ↓
                           PostgreSQL INSERT
\end{verbatim}

\section{Fluxo de Dados - Comando}

\begin{verbatim}
Dashboard Web ─── POST /api/command ──→ Backend
                                          ↓
                        Publish via MQTT ao tópico de comando
                                          ↓
                                    Mosquitto
                                          ↓
                        ESP32 subscrito em /command recebe
\end{verbatim}

\section{Fluxo de Dados - Autenticação}

\begin{verbatim}
Usuário ─── POST /login ──→ Backend (valida em memória)
                              ↓
                         Set sessão (cookie/token)
                              ↓
                   Redirecion para dashboard
\end{verbatim}

% ============================================================================
% CAPÍTULO 7
% ============================================================================
\chapter{Mapeamento de Conceitos de Sistemas Distribuídos}

\begin{center}
\begin{tabular}{|m{3.5cm}|m{9.5cm}|}
    \hline
    \rowcolor{lightblue}
    \textbf{Conceito} & \textbf{Implementação} \\
    \hline
    Processos/Threads & Node.js (event-loop single-thread, cluster para múltiplas instâncias) \\
    \hline
    Sockets TCP/UDP & MQTT (TCP), HTTP (TCP) \\
    \hline
    RPC & HTTP REST (futuro: gRPC) \\
    \hline
    Message Queue & MQTT broker (Mosquitto) \\
    \hline
    Naming Service & DNS + Tailscale MagicDNS (service discovery) \\
    \hline
    Sincronização & Shared state em memória + BD (eventual consistency) \\
    \hline
    Exclusão Mútua & PostgreSQL locks, transações ACID \\
    \hline
    Replicação & Futura: PostgreSQL replication + Mosquitto clustering \\
    \hline
    Tolerância a Falhas & Health-checks, auto-restart, reconexão \\
    \hline
    Segurança & RBAC (papéis), sessões, Tailscale (VPN) \\
    \hline
\end{tabular}
\end{center}

% ============================================================================
% CAPÍTULO 8
% ============================================================================
\chapter{Tecnologias Utilizadas}

\section{Stack Tecnológico}

\begin{itemize}
    \item \textbf{Linguagem de Programação:} JavaScript (Node.js v18+)
    \item \textbf{Banco de Dados:} PostgreSQL 16 + TimescaleDB (série temporal)
    \item \textbf{Message Broker:} Mosquitto (MQTT v3.1.1)
    \item \textbf{Firmware Embarcado:} Arduino/PlatformIO (ESP32)
    \item \textbf{Frontend:} HTML5 + Vanilla JavaScript
    \item \textbf{Acesso Remoto:} Tailscale (Zero-Trust VPN)
    \item \textbf{Orquestração:} Docker Compose (futura) / Systemd (atual)
\end{itemize}

\section{Infraestrutura de Rede}

\begin{itemize}
    \item \textbf{Rede Local:} WiFi 192.168.1.0/24
    \item \textbf{Rede Remota:} Tailscale VPN com 100.82.x.x/16
    \item \textbf{Acesso Público:} Tailscale Funnel (HTTPS)
\end{itemize}

% ============================================================================
% CAPÍTULO 9
% ============================================================================
\chapter{Próximas Etapas Previstas}

\section{Etapa 2 (Concorrência e Sockets)}

Implementação de:

\begin{itemize}
    \item Múltiplos workers/threads para processamento paralelo de telemetria
    \item Pool de conexões TCP/MQTT com timeout e reconnect
    \item Load balancing entre múltiplos backends (futuro: nginx)
\end{itemize}

\section{Etapa 3 (Sincronização e RPC)}

\begin{itemize}
    \item Consensus entre múltiplas instâncias do backend (Raft/Paxos)
    \item RPC tipado com gRPC
    \item Transações distribuídas (two-phase commit)
\end{itemize}

\section{Etapa 4 (Replicação e Cache)}

\begin{itemize}
    \item Replicação de BD entre Raspberry Pi 3 e máquina backup
    \item Cache distribuído com Redis
    \item CDN para assets estáticos
\end{itemize}

\section{Etapa 5 (Tolerância a Falhas)}

\begin{itemize}
    \item Circuit breaker, bulkhead, retry patterns
    \item Health-check automático e failover
    \item Chaos engineering para testes de resiliência
\end{itemize}

\section{Etapa 6 (Segurança Avançada)}

\begin{itemize}
    \item TLS/SSL para MQTT e HTTP
    \item OAuth2 / OpenID Connect
    \item Auditoria e logging centralizado
    \item Rate limiting e DDoS mitigation
\end{itemize}

% ============================================================================
\newpage
\appendix

\chapter{Referências Técnicas}

\begin{itemize}
    \item \textbf{MQTT Specification:} \url{https://mqtt.org/}
    \item \textbf{PostgreSQL Documentation:} \url{https://www.postgresql.org/docs/}
    \item \textbf{TimescaleDB Documentation:} \url{https://docs.timescale.com/}
    \item \textbf{Node.js API:} \url{https://nodejs.org/api/}
    \item \textbf{Tailscale Documentation:} \url{https://tailscale.com/kb/}
    \item \textbf{ESP32 PlatformIO:} \url{https://docs.platformio.org/}
\end{itemize}

\chapter{Configuração de Infraestrutura}

\section{Máquinas Envolvidas}

\begin{center}
\begin{tabular}{|l|l|l|l|}
    \hline
    \rowcolor{lightblue}
    \textbf{Máquina} & \textbf{IP Local} & \textbf{IP Tailscale} & \textbf{Função} \\
    \hline
    Notebook & 192.168.1.X & 100.82.199.70 & Web + Backend \\
    \hline
    Raspberry Pi 3 & 192.168.1.10 & 100.82.140.119 & Banco de Dados \\
    \hline
    ESP32 Node & 192.168.1.Y & - & Sensores \\
    \hline
\end{tabular}
\end{center}

\section{Portas Utilizadas}

\begin{center}
\begin{tabular}{|l|l|l|}
    \hline
    \rowcolor{lightblue}
    \textbf{Serviço} & \textbf{Porta} & \textbf{Protocolo} \\
    \hline
    Dashboard Web & 3000 & HTTP \\
    \hline
    API Gateway & 3001 & HTTP \\
    \hline
    MQTT Broker & 1883 & MQTT/TCP \\
    \hline
    PostgreSQL & 5432 & PostgreSQL/TCP \\
    \hline
\end{tabular}
\end{center}

\end{document}
